\documentclass[10pt]{article}
\usepackage[margin=0.55in]{geometry}
\usepackage{amsmath,amssymb}
\usepackage[T1]{fontenc}
\usepackage[hidelinks]{hyperref}
\usepackage{microtype}

\newcommand{\M}{\mathcal M}

\title{A Genuine Noncommutative Hartman--Wintner LIL\\
Was Obtained in 2026}
\author{fa\_banach\_001 / agent\_lane\_00}
\date{17 August 2026}

\begin{document}
\maketitle
\vspace{-1.2em}
\small

\noindent\textbf{Status: literature already answered (affirmative answer
to the source's existence question; no new theorem is claimed here).}

\section*{Original question}

Section~4 of Zeng's paper asks whether a ``genuine'' noncommutative
Hartman--Wintner law of the iterated logarithm is possible, where genuine
means that the theorem does not assume Kolmogorov's pointwise small-increment
condition
\[
                   \|d_n\|_\infty\leq \alpha_n s_n/u_n,
                   \qquad \alpha_n\longrightarrow0.
\]
The source points to Stout's classical theorem for stationary-ergodic
martingale differences and leaves the noncommutative possibility open
\cite{Zeng} (source PDF page~8).

\section*{Explicit later answer}

Panja, Ricard, and Saha explicitly cite Zeng and prove optimal
noncommutative LILs in arXiv:2509.22037v2 \cite{PRS}.

First, their Theorem~3.4 (supporting PDF page~12) is a genuine martingale
version obtained by cutting large differences.  Instead of the pointwise
small-increment assumption, it imposes a conditional Lindeberg-type tail
summability condition: for some $K_n\to0$ and $\alpha>0$,
\[
 \left\|\sum_{n\geq1}\frac{u_n^\alpha}{s_n^2}
 E_{n-1}\!\left[d_n^2
 \mathbf 1_{(s_nK_n/u_n,\infty)}(|d_n|)\right]\right\|<\infty.
\]
It concludes the optimal almost-uniform upper bound
\[
          \limsup_{n\to\infty}\frac{x_n}{s_nu_n}
          \underset{\mathrm{a.u.}}{\leq}\sqrt2.
\]

Second, Theorem~4.4 (supporting PDF page~14) removes the small-increment
condition completely for independent identically distributed data.  If
$(y_i)$ are independent, identically distributed, mean-zero
noncommutative random variables with $\|y_i\|_2\leq1$, then, writing
$x_n=\sum_{i=1}^n y_i$ and $u_n=\sqrt{\max\{1,\log\log n\}}$,
\[
          \limsup_{n\to\infty}\frac{x_n}{\sqrt n\,u_n}
          \underset{\mathrm{a.u.}}{\leq}\sqrt2.
\]
Independent identically distributed differences are stationary and
ergodic.  Thus the source's question whether a genuine noncommutative
Hartman--Wintner theorem is \emph{possible} has an affirmative answer.

\section*{Scope}

The later paper proves an optimal genuine martingale theorem under a tail
condition and the full IID $L_2$ Hartman--Wintner theorem.  It does not claim
the strongest conceivable statement for every dependent stationary-ergodic
noncommutative martingale-difference sequence.  The source's open-ended
possibility question is answered, while that broader universal extension
should not be inferred.

The later authors also identify a slight gap in the exponential-inequality
application used for Zeng's original constant-$2$ result; their refined
argument both corrects the route and obtains the classical sharp constant.

\scriptsize
\begin{thebibliography}{9}
\bibitem{Zeng}
Q. Zeng,
\emph{Kolmogorov's law of the iterated logarithm for noncommutative
martingales}, Ann. Inst. H. Poincar\'e Probab. Statist. \textbf{51} (2015),
1124--1130; arXiv:1212.1504v3.

\bibitem{PRS}
S. Panja, \'E. Ricard, and D. Saha,
\emph{Non-commutative Law of iterated logarithm},
J. Funct. Anal. (2026), article 111477; arXiv:2509.22037v2.
\end{thebibliography}

\end{document}
